\chapter{Triết Lý Sống (Philosophy of Life)}

\begin{center}
  \textit{\color{truyenblue}``The unexamined life is not worth living.''}\\
  \textit{(Một cuộc đời không được soi xét thì không đáng sống.)}\\[4pt]
  \small\color{truyengold}--- Socrates ---
\end{center}

\ngancach

\section{The Cup of Tea (Tách Trà)}
\begin{truyen}{Triết Lý Sống}{Thiết Thực}
A wise man \textbf{pours} tea into a \textbf{cup} until it \textbf{overflows}. He says you must \textbf{empty} your mind to \textbf{learn} new things.

\textit{Một người thông thái rót trà vào tách cho đến khi tràn ra ngoài. Ông nói rằng bạn phải làm trống tâm trí của mình để học hỏi những điều mới.}

\begin{baihoc}
Hãy luôn giữ tinh thần học hỏi với một tâm trí cởi mở.
\end{baihoc}
\end{truyen}

\section{The Two Wolves (Hai Con Sói)}
\begin{truyen}{Triết Lý Sống}{Lựa Chọn}
An old man tells his \textbf{grandson} about two \textbf{wolves} fighting inside us: one is \textbf{good} and one is \textbf{bad}. The one you \textbf{feed} will \textbf{win}.

\textit{Một ông lão kể cho cháu trai nghe về hai con sói đang chiến đấu bên trong chúng ta: một con tốt và một con xấu. Con nào bạn nuôi dưỡng sẽ chiến thắng.}

\begin{baihoc}
Những gì bạn nuôi dưỡng trong tâm hồn sẽ định hình con người bạn.
\end{baihoc}
\end{truyen}

\section{The Heavy Glass (Ly Nước Nặng)}
\begin{truyen}{Triết Lý Sống}{Buông Bỏ}
A teacher \textbf{holds} a \textbf{glass} of water. She says the \textbf{weight} depends on how \textbf{long} you hold it.

\textit{Một giáo viên cầm một ly nước. Cô nói rằng trọng lượng phụ thuộc vào việc bạn cầm nó bao lâu.}

\begin{baihoc}
Đừng giữ những muộn phiền quá lâu, hãy học cách buông bỏ để thấy nhẹ nhõm hơn.
\end{baihoc}
\end{truyen}

\section{The Mirror in the Room (Tấm Gương)}
\begin{truyen}{Triết Lý Sống}{Phán Xét}
A man \textbf{looks} in a \textbf{mirror} and sees his \textbf{flaws}. He realizes he must \textbf{fix} himself before \textbf{judging} others.

\textit{Một người đàn ông nhìn vào gương và thấy những khuyết điểm của mình. Anh nhận ra mình phải sửa bản thân trước khi phán xét người khác.}

\begin{baihoc}
Hãy nhìn nhận lại chính mình trước khi chỉ trích người khác.
\end{baihoc}
\end{truyen}

\section{The Broken Vase (Chiếc Bình Vỡ)}
\begin{truyen}{Triết Lý Sống}{Tổn Thương}
A \textbf{beautiful} vase is \textbf{broken}. An artist \textbf{repairs} it with \textbf{gold}, making it more \textbf{valuable} than before.

\textit{Một chiếc bình đẹp bị vỡ. Một nghệ nhân sửa nó bằng vàng, làm cho nó có giá trị hơn trước.}

\begin{baihoc}
Những khó khăn và tổn thương có thể làm cho chúng ta trở nên mạnh mẽ và đáng giá hơn.
\end{baihoc}
\end{truyen}

\section{The Blind Men and the Elephant (Thầy Bói Xem Voi)}
\begin{truyen}{Triết Lý Sống}{Góc Nhìn}
Six \textbf{blind} men \textbf{touch} an elephant. Each thinks it is \textbf{something} else. They \textbf{need} all parts to know the \textbf{truth}.

\textit{Sáu người mù sờ một con voi. Mỗi người nghĩ nó là một thứ khác. Họ cần tất cả các phần để biết sự thật.}

\begin{baihoc}
Mỗi người đều có một góc nhìn riêng, hãy lắng nghe để hiểu được bức tranh toàn cảnh.
\end{baihoc}
\end{truyen}

\section{The Footprints in the Sand (Dấu Chân Trên Cát)}
\begin{truyen}{Triết Lý Sống}{Khó Khăn}
A man \textbf{walks} on the \textbf{beach} and looks back. He sees that in \textbf{hard} times, there is only \textbf{one} set of footprints.

\textit{Một người đàn ông đi dạo trên bãi biển và nhìn lại. Anh thấy rằng trong những lúc khó khăn, chỉ có một hàng dấu chân.}

\begin{baihoc}
Trong những lúc khó khăn nhất, đôi khi bạn cần tự bước đi bằng chính đôi chân của mình.
\end{baihoc}
\end{truyen}

\section{The Empty Boat (Chiếc Thuyền Trống)}
\begin{truyen}{Triết Lý Sống}{Bình Yên}
A man gets \textbf{angry} when a \textbf{boat} hits his. When he sees it is \textbf{empty}, his \textbf{anger} goes \textbf{away}.

\textit{Một người đàn ông tức giận khi một chiếc thuyền đâm vào thuyền của mình. Khi anh thấy nó trống rỗng, cơn giận của anh biến mất.}

\begin{baihoc}
Đừng lãng phí năng lượng vào sự vung vãi vô tình của cuộc đời, hãy giữ tâm hồn bình yên.
\end{baihoc}
\end{truyen}

\section{The Three Questions (Ba Câu Hỏi)}
\begin{truyen}{Triết Lý Sống}{Hiện Tại}
A king \textbf{asks} three \textbf{questions} to find happiness. He \textbf{learns} the most \textbf{important} time is \textbf{now}.

\textit{Một vị vua hỏi ba câu hỏi để tìm hạnh phúc. Ông học được rằng thời điểm quan trọng nhất là hiện tại.}

\begin{baihoc}
Hãy sống trọn vẹn ở hiện tại để tìm thấy hạnh phúc thực sự.
\end{baihoc}
\end{truyen}

\section{The River and the Rock (Sông Và Đá)}
\begin{truyen}{Triết Lý Sống}{Kiên Trì}
A \textbf{river} hits a \textbf{rock}. Over \textbf{time}, the water \textbf{cuts} through the rock with \textbf{persistence}.

\textit{Một con sông đập vào một hòn đá. Qua thời gian, dòng nước xuyên qua hòn đá bằng sự kiên trì.}

\begin{baihoc}
Sự kiên trì và bền bỉ có thể vượt qua mọi chướng ngại vật cứng rắn nhất.
\end{baihoc}
\end{truyen}
